\scanpage{45}
\noindent\hfill atkin \mspage{5}

Newman proves \handwritten{(3.2)} by observing that $t(n)\handwritten{\equiv\tau}(n)\;(\mathrm{mod}\;13)$, where* $\handwritten{\tau}(n)$ is defined by

\noindent\handwritten{\rule{\linewidth}{0.4pt}}

\noindent *We apologise to the reader for the two uses of $\handwritten{\tau}$. Both are well established, and there is no real danger of confusion.

\noindent\handwritten{\rule{\linewidth}{0.4pt}}

\[
\handwritten{\sum_{n=1}^{\infty}\tau(n)\,x^{n}=xf^{24}(x)}
\]

Now $\handwritten{\tau}(n)$ s\rlap{\handwritten{\tikz[baseline=0pt,x=1em,y=1ex]{\draw (0,-.35)--(.56,-.35)--(.56,1.7)--(1.06,1.7);}}}taisfies the \uline{identical} relation
\[
\handwritten{\tau(np)-\tau(n)\tau(p)+p^{11}\tau(n/p)=0}
\tag*{\handwritten{(3.5)}}
\]
as conjectured by RAMANUJAN(1916) and proved by MORDELL(1919),so that \handwritten{(3.1)} and \handwritten{(3.2)} are immediate. We may remark here that Ramanujan probably found this conjecture also by inspection of his own table of \deleted{e}\handwritten{\unsure{v|V}}alues of $\handwritten{\tau}(n)$ and his extraordinary intuition for algebraic form in this subject. It is a consequence of \handwritten{(3.5)} that $\handwritten{\tau}(n)$ is multiplicative,i.e.,
\[
\handwritten{\tau(mn)=\tau(m)\tau(n)\quad\text{if}\quad (m,n)=1.}
\tag*{\handwritten{(3.6)}}
\]

GLAISHER(\handwritten{1907}) had previously obtained a number of results similar to \handwritten{(3.6)} for coefficients relating to the expression of a number as sums of squares,but he did not obtain the crucial third term ( analogous to $p^{11}\handwritten{\tau}(n/p)$ in \handwritten{(3.5)})in his equations to obtain the full result similar to \handwritten{(3.5)}. Glaisher's results have been completed and fully proved by RANKIN(196\hspace{1em}).

% archive sheet number: 44
